\documentclass[11pt,a4paper]{article}

\usepackage[margin=2.5cm]{geometry}
\usepackage{graphicx}
\usepackage{booktabs}
\usepackage{amsmath,amssymb}
\usepackage{hyperref}
\usepackage{float}
\usepackage{subcaption}
\usepackage{xcolor}
\usepackage{siunitx}
\usepackage{enumitem}
\usepackage{attachfile2}

\hypersetup{colorlinks=true, linkcolor=blue!60!black, citecolor=blue!60!black, urlcolor=blue!60!black}

\title{Physically Grounded Flexible Cable Simulation\\in NVIDIA Isaac Sim}
\author{Ali Bani Asad\\[4pt]
\small Prospective PhD Student, University of Waterloo\\
\small Supervisor: Prof.\ Soo Jeon}
\date{May 2026}

\begin{document}

\maketitle

\begin{abstract}
This report presents a method for simulating flexible cables with rigid connectors
in NVIDIA Isaac Sim using chains of rigid capsules connected by D6 joints whose
stiffness and damping are derived from material properties via Euler--Bernoulli
beam theory. I reproduce the parameter sweep of Govoni et al.\ (2025) and show
that the PhysX TGS solver remains stable in all five tested configurations,
including one regime where the reference paper's explicit mass-spring-damper
integrator diverges. Beyond the reproduction study, I present a parametric
analysis of the derived joint parameters across material classes, dynamic
trajectory analysis of cable behaviour, and a damping sensitivity study
that quantifies the error in hand-tuned approaches. These results establish
a solid foundation for future reinforcement-learning-based cable manipulation
research.
\end{abstract}


% ------------------------------------------------------------------
\section{Introduction}
% ------------------------------------------------------------------

Simulating deformable linear objects (DLOs) such as cables, wires, and hoses
is a long-standing challenge in robotics simulation.  NVIDIA Isaac Sim does not
natively support one-dimensional deformable bodies, so cables must be
approximated using chains of rigid segments connected by constrained joints.

This work develops such a model with three key goals:
\begin{enumerate}[nosep]
    \item \textbf{Physical grounding} --- joint parameters are derived from
          measurable material properties (Young's modulus, geometry, damping
          ratio), not hand-tuned.
    \item \textbf{Stability validation} --- the model is benchmarked against the
          systematic parameter sweep of Govoni et al.~\cite{govoni2025},
          which documents stability boundaries for mass-spring-damper (MSD)
          cable models~\cite{lv2017} across different stiffness regimes.
    \item \textbf{Behavioural characterisation} --- dynamic cable behaviour is
          analysed through trajectory data, shape evolution, and energy
          dissipation to verify physical plausibility beyond mere stability.
\end{enumerate}


% ------------------------------------------------------------------
\section{Modelling Approach}
% ------------------------------------------------------------------

\subsection{Capsule-Chain Architecture}

The cable is modelled as a chain of $N$ capsule-shaped rigid bodies connected
by \textbf{D6 joints} (six degree-of-freedom joints).  Each joint locks all
three translational axes, enforcing inextensibility, and constrains bending
via a soft PD drive plus a hard cone limit as a safety net.

Two rigid connectors are attached at the cable endpoints:
\begin{itemize}[nosep]
    \item \textbf{Top connector} (blue, \SI{3}{cm} cube, \SI{2}{kg}): fixed to the world frame.
    \item \textbf{Bottom connector} (red, \SI{2}{cm} cube, \SI{0.2}{kg}): free
          or kinematic, depending on experiment mode.
\end{itemize}

Following Govoni et al., \textbf{twist degrees of freedom are left free}
(no twist limit, no rotational-X drive), since twisting springs are unnecessary
for typical DLO manipulation tasks.

\subsection{Material-Driven Parameter Derivation}

Rather than hand-picking joint stiffness and damping values, I derive them
from material properties using Euler--Bernoulli beam theory.  Alternative
discrete-rod formulations~\cite{bergou2008} derive bending energies from
material-frame curvatures rather than joint angles.  Given:
\begin{itemize}[nosep]
    \item Young's modulus $E$ (\si{Pa}),
    \item Cable radius $r$ (\si{m}),
    \item Segment length $L_{\text{seg}} = L_{\text{total}} / N$ (\si{m}),
    \item Damping ratio $\zeta$ (fraction of critical damping),
\end{itemize}
the joint parameters are computed as:

\begin{align}
    I_{\text{area}} &= \frac{\pi r^4}{4}
        & &\text{(second moment of area)} \label{eq:I}\\[4pt]
    K_{\text{bend}} &= \frac{E \cdot I_{\text{area}}}{L_{\text{seg}}}
        & &\text{(bending stiffness, \si{N.m/rad})} \label{eq:K}\\[4pt]
    I_{\text{rot}} &= \tfrac{1}{3}\, m_{\text{seg}}\, L_{\text{seg}}^2
        & &\text{(rotational inertia of segment)} \label{eq:Irot}\\[4pt]
    C_{\text{crit}} &= 2\sqrt{K_{\text{bend}} \cdot I_{\text{rot}}}
        & &\text{(critical damping~\cite{bhasin2006})} \label{eq:Ccrit}\\[4pt]
    C &= \zeta \cdot C_{\text{crit}}
        & &\text{(joint damping, \si{N.m.s/rad})} \label{eq:C}
\end{align}

Both $K_{\text{bend}}$ and $C$ are converted from \si{rad} to \si{deg}
(multiply by $\pi/180$) before being passed to the USD DriveAPI, since
Isaac Sim uses degrees for all angular quantities.


\subsection{Comparison with v1 Baseline}

My initial simulation (v1) used hand-picked parameters: zero bending
stiffness ($K = 0$), a damping of \SI{0.05}{N.m.s/deg}, a narrow
\ang{8} cone limit, and a \ang{5.8} twist limit.  Analysis revealed that
the damping value was approximately $70{,}000\times$ critical damping
for the cable's geometry and material --- the cable behaved as if submerged
in viscous fluid rather than air.

Figure~\ref{fig:damping_sensitivity} quantifies this discrepancy.  The
physically correct damping for the rubber cable at $\zeta = 0.2$ is
\SI{1.43e-7}{N.m.s/deg}, while the v1 value of \SI{0.05}{N.m.s/deg} is
orders of magnitude above even critical damping ($\zeta = 1$).

\begin{figure}[H]
\centering
\includegraphics[width=0.85\textwidth]{figures/damping_sensitivity.pdf}
\caption{Joint damping coefficient as a function of damping ratio for
the default rubber cable ($E = \SI{12.6}{MPa}$, $N = 200$).  The v1
hand-tuned value (red dotted line) exceeds critical damping by a factor
of $\sim\!70{,}000$, producing unrealistic overdamped behaviour.}
\label{fig:damping_sensitivity}
\end{figure}


% ------------------------------------------------------------------
\section{Parametric Analysis}
% ------------------------------------------------------------------

Before running simulations, I analyse how the derived joint parameters
vary across the space of materials and discretisations.  This analysis
informs which configurations are feasible at a given physics timestep
and guides domain-randomisation ranges for future RL training.

\begin{figure}[H]
\centering
\includegraphics[width=\textwidth]{figures/stiffness_landscape.pdf}
\caption{\textbf{Left:} Derived bending stiffness $K_{\text{bend}}$ as a
function of Young's modulus for different discretisations $N$.  Black markers
show the three Govoni test materials.  Shaded regions indicate material
classes.  \textbf{Right:} Natural frequency $f_n = \frac{1}{2\pi}\sqrt{K/I_{\text{rot}}}$
of each joint segment.  The dashed line marks the Nyquist limit of the
physics solver ($240/2 = 120$~Hz); configurations above this line risk
under-sampling the joint dynamics.}
\label{fig:stiffness_landscape}
\end{figure}

Key observations from Figure~\ref{fig:stiffness_landscape}:
\begin{itemize}[nosep]
    \item For a fixed material, finer discretisation (higher $N$) reduces
          per-joint stiffness (shorter segments are less stiff), but increases
          computational cost.
    \item The Govoni test points (black markers) span three orders of magnitude
          in stiffness, from soft rubber to rigid thermoplastics.
    \item At $N = 200$, the natural frequency remains below the Nyquist limit
          for $E$ up to approximately \SI{25}{MPa} (rubber regime); stiffer plastic
          cables exceed Nyquist at $N = 200$ and would benefit from coarser
          discretisation.
    \item The high-stiffness regime requires explicit integrators to use
          timesteps $\Delta t \propto 1/\sqrt{k}$, which is exactly what
          Govoni's Row~3 vs.\ Row~4 result demonstrates: stability is
          recovered only by reducing $\Delta t$ from \SI{5e-6}{s} to
          \SI{1e-7}{s}.
\end{itemize}


% ------------------------------------------------------------------
\section{Experiment Design}
% ------------------------------------------------------------------

\subsection{Experiment Modes}

The simulation supports two modes:
\begin{description}[nosep]
    \item[Hanging kick] The top connector is fixed to the world; the bottom
        connector is free and receives an initial lateral velocity of
        \SI{1.5}{m/s}.  A green obstacle cube is placed in the scene to test
        contact interaction.  This mode demonstrates visual realism and
        dynamic behaviour.
    \item[Both ends fixed (Govoni-style)] Both connectors are fixed.  The cable
        settles under gravity for \SI{0.5}{s}, then a \SI{5}{mm} step
        displacement is applied to the bottom connector.  This replicates the
        stability test protocol of Govoni et al.
\end{description}

\subsection{Govoni Sweep Reproduction}

To validate the model, I reproduce the five configurations from
Table~1 of Govoni et al.~\cite{govoni2025}.  Each configuration varies
Young's modulus $E$ and the number of discrete masses $i$:

\begin{table}[H]
\centering
\caption{Govoni et al.\ Table 1 configurations.}
\label{tab:govoni_configs}
\begin{tabular}{@{}clcc@{}}
\toprule
\textbf{Row} & \textbf{$E$ (MPa)} & \textbf{$i$ (links)} & \textbf{Paper result} \\
\midrule
1  & 12.6   & 10 & Stable \\
2a & 526.0  & 10 & Stable \\
2b & 1002.0 &  6 & Stable \\
3  & 1002.6 & 10 & \textcolor{red}{\textbf{Unstable}} (at \SI{0.4}{s}) \\
4  & 1002.6 & 10 & Stable \\
\bottomrule
\end{tabular}
\end{table}

Row~3 is the critical case: Govoni's explicit MSD integrator diverges after
\SI{0.4}{s} at $\Delta t = \SI{5e-6}{s}$.  Row~4 recovers stability by
reducing $\Delta t$ to \SI{1e-7}{s}.

The sweep uses the PhysX TGS solver at $\Delta t = 1/240\;\text{s}
\approx \SI{4.17e-3}{s}$, with 64 position iterations and 8 velocity
iterations.  An automated driver script (\texttt{govoni\_sweep.py}) launches
each configuration, monitors angular velocities for divergence, and assembles
a comparison table.


% ------------------------------------------------------------------
\section{Results}
% ------------------------------------------------------------------

\subsection{Govoni Sweep: Stability Comparison}

\begin{table}[H]
\centering
\caption{Reproduction of Govoni et al.\ Table 1.  The proposed solver operates at
$\Delta t = 1/240\;\text{s}$ (PhysX TGS, 64 position iterations), compared
to their explicit MSD at $\Delta t = \SI{5e-6}{s}$.}
\label{tab:comparison}
\begin{tabular}{@{}clccccc@{}}
\toprule
\textbf{Row} & \textbf{$E$ (MPa)} & \textbf{$i$}
    & \textbf{Paper} & \textbf{Mine}
    & \textbf{Max $|\omega|$ (deg/s)}
    & \textbf{Match} \\
\midrule
1  & 12.6   & 10 & Stable   & Stable & $5.34 \times 10^2$ & \checkmark \\
2a & 526.0  & 10 & Stable   & Stable & $3.39 \times 10^2$ & \checkmark \\
2b & 1002.0 &  6 & Stable   & Stable & $1.78 \times 10^2$ & \checkmark \\
3  & 1002.6 & 10 & \textcolor{red}{Unstable} & \textcolor{green!60!black}{\textbf{Stable}}
    & $2.57 \times 10^2$ & $\star$ \\
4  & 1002.6 & 10 & Stable   & Stable & $2.57 \times 10^2$ & \checkmark \\
\bottomrule
\end{tabular}
\end{table}

\noindent
All five configurations remain stable under my solver.  Rows 1, 2a, 2b, and
4 reproduce the paper's findings directly.

\textbf{Row~3 is the key result}: where Govoni's explicit integrator diverges
after \SI{0.4}{s}, my PhysX TGS solver remains stable with a peak angular
velocity of only \SI{257}{deg/s} --- well within normal operating range.  This
is expected: PhysX TGS is an \emph{implicit} iterative solver that naturally
handles stiff systems, whereas an explicit integrator requires
$\Delta t \propto 1/\sqrt{k}$ for stability.

Figure~\ref{fig:max_omega} visualises the stability margin across all
configurations.

\begin{figure}[H]
\centering
\includegraphics[width=0.85\textwidth]{figures/max_omega_bar.pdf}
\caption{Peak angular velocity recorded during each Govoni sweep
configuration.  All values remain far below the divergence threshold
(\SI{e4}{deg/s}).  Row~3 (red border) is the configuration where
Govoni's explicit integrator diverges, yet my solver produces peak
velocities comparable to the stable configurations.}
\label{fig:max_omega}
\end{figure}


\subsection{Govoni Sweep: Dynamic Response}

Beyond binary stability, I examine the transient response of each
configuration after the \SI{5}{mm} step displacement.

\begin{figure}[H]
\centering
\includegraphics[width=0.85\textwidth]{figures/govoni_midpoint.pdf}
\caption{Midpoint vertical displacement relative to initial position for all
five Govoni configurations.  Stiffer cables (higher $E$) respond faster but
with smaller amplitude.  Row~3 (dashed red) and Row~4 (dotted purple) are
identical in my solver since both use the same material; the paper's
instability at Row~3 was a solver limitation, not a physical phenomenon.}
\label{fig:govoni_midpoint}
\end{figure}


\subsection{Hanging Cable: Trajectory Analysis}

The 200-link hanging-kick experiment ($E = \SI{12.6}{MPa}$, rubber) provides
rich dynamic data.  Figure~\ref{fig:keyframes} shows key frames from the
\SI{10}{s} simulation.

\begin{figure}[H]
\centering
\begin{subfigure}[b]{0.48\textwidth}
    \includegraphics[width=\textwidth]{../cable_output/hanging_kick/frame_t2s.png}
    \caption{$t = \SI{2}{s}$ --- cable swinging after kick}
\end{subfigure}
\hfill
\begin{subfigure}[b]{0.48\textwidth}
    \includegraphics[width=\textwidth]{../cable_output/hanging_kick/frame_t5s.png}
    \caption{$t = \SI{5}{s}$ --- contact with obstacle}
\end{subfigure}

\vspace{0.5cm}

\begin{subfigure}[b]{0.48\textwidth}
    \includegraphics[width=\textwidth]{../cable_output/hanging_kick/frame_t10s.png}
    \caption{$t = \SI{10}{s}$ --- settling to equilibrium}
\end{subfigure}
\caption{Key frames from the 200-link hanging cable simulation
($E = \SI{12.6}{MPa}$, $\zeta = 0.2$, \SI{10}{s} run).  The simulation
remained stable throughout, with $\max|\omega| = \SI{2093}{deg/s}$.}
\label{fig:keyframes}
\end{figure}

A full \SI{10}{s} video recording of the simulation is attached to this PDF
(\textattachfile[color=0 0.4 0.8, description={10s cable simulation video},
mimetype=video/mp4, print=false]{../cable_output/hanging_kick/cable_simulation.mp4}%
{\textbf{cable\_simulation.mp4} --- click to open}).

Figure~\ref{fig:tip_trajectory} presents the quantitative trajectory data
extracted from the simulation log.

\begin{figure}[H]
\centering
\includegraphics[width=0.85\textwidth]{figures/tip_trajectory.pdf}
\caption{Cable tip and midpoint positions over \SI{10}{s}.  \textbf{Top:}
lateral (X) displacement shows the cable swinging in its pendulum mode, which
is damped primarily by translational drag (LINEAR\_DAMPING in PhysX) rather
than by joint stiffness; per-joint angular damping dissipates bending energy
but couples weakly to bulk pendulum motion.  \textbf{Bottom:} vertical (Z)
positions show the tip oscillating around its vertical rest position with small
amplitude, with higher-frequency bending modes well-damped at $\zeta = 0.2$.}
\label{fig:tip_trajectory}
\end{figure}


\subsection{Cable Shape Evolution}

Figure~\ref{fig:cable_shapes} shows the cable's spatial configuration at
selected time instants, using the logged positions of tracked capsules.

\begin{figure}[H]
\centering
\includegraphics[width=0.75\textwidth]{figures/cable_shapes.pdf}
\caption{Cable shape at selected time instants during the hanging-kick
experiment.  The cable transitions from a vertical rest position through
large lateral swings before settling into a gravity-equilibrium shape.}
\label{fig:cable_shapes}
\end{figure}


\subsection{Phase Portrait: Energy Dissipation}

To visualise the energy dissipation behaviour, Figure~\ref{fig:phase_portrait}
shows the cable tip trajectory in the X--Z plane, colour-coded by time.

\begin{figure}[H]
\centering
\includegraphics[width=0.7\textwidth]{figures/phase_portrait.pdf}
\caption{Cable tip trajectory in the X--Z plane, colour-coded by time.
The pattern shows the tip oscillating across the swing plane with gradual
concentration toward equilibrium near the end of the simulation, indicating
partial energy dissipation.  The lack of a clean spiral reflects the
multi-modal nature of the cable's motion: pendulum swing, bending modes,
and gravity-driven re-equilibration overlap in this projection.}
\label{fig:phase_portrait}
\end{figure}


% ------------------------------------------------------------------
\section{Simulation Parameters}
% ------------------------------------------------------------------

\begin{table}[H]
\centering
\caption{Cable and solver parameters (material-driven).}
\label{tab:params}
\begin{tabular}{@{}llr@{}}
\toprule
\textbf{Parameter}         & \textbf{Symbol / Note}    & \textbf{Default Value}  \\
\midrule
\multicolumn{3}{@{}l}{\textit{Cable geometry}} \\
Number of links            &  $N$               & 200             \\
Total cable length         &  $L$               & \SI{1.0}{m}     \\
Link radius                &  $r$               & \SI{1.5}{mm}    \\
Segment spacing            &  $L_{\text{seg}}$  & \SI{5.0}{mm}    \\
Total cable mass           &  $m_{\text{total}}$ & \SI{1.0}{kg}   \\
\midrule
\multicolumn{3}{@{}l}{\textit{Material properties}} \\
Young's modulus            &  $E$               & \SI{12.6}{MPa} (rubber) \\
Poisson's ratio            &  $\nu$             & 0.5             \\
Damping ratio              &  $\zeta$           & 0.2             \\
\midrule
\multicolumn{3}{@{}l}{\textit{Derived joint parameters (from Eqs.~\ref{eq:I}--\ref{eq:C})}} \\
Bending stiffness          &  $K_{\text{bend}}$ & \SI{1.75e-4}{N.m/deg}   \\
Joint damping              &  $C$               & \SI{1.43e-7}{N.m.s/deg} \\
Cone limit (safety net)    &  $\theta_c$        & \ang{30}        \\
Twist                      &  ---               & Free (no limit) \\
\midrule
\multicolumn{3}{@{}l}{\textit{PhysX solver}} \\
Solver type                &                    & TGS             \\
Position iterations        &                    & 64              \\
Velocity iterations        &                    & 8               \\
Physics timestep           &  $\Delta t$        & $1/240$ s       \\
CCD                        &                    & Enabled         \\
\bottomrule
\end{tabular}
\end{table}


% ------------------------------------------------------------------
\section{Discussion}
% ------------------------------------------------------------------

\subsection{Advantages of the Implicit Solver Approach}

The explicit MSD formulation used by Govoni et al.\ requires timesteps on the
order of microseconds to remain stable with stiff materials.  At
$\Delta t = \SI{5e-6}{s}$, simulating \SI{1}{s} of cable behaviour requires
200{,}000 integration steps.  My approach leverages PhysX's implicit TGS
solver, which remains stable at $\Delta t \approx \SI{4e-3}{s}$ ---
approximately $800\times$ larger.  For Row~4, where Govoni reduces
$\Delta t$ to $\SI{1e-7}{s}$ to recover stability, my solver runs at a
timestep approximately $41{,}000\times$ larger.  This makes real-time or
faster-than-real-time simulation feasible, which is essential for
reinforcement learning where millions of timesteps of experience are needed.

\subsection{Implications for Domain Randomisation}

A key concern raised by Govoni et al.\ is that domain randomisation over
material properties (varying $E$ during RL training) can produce unstable
simulation samples that silently corrupt the training set.  My results
show that this is a solver-specific problem, not a fundamental limitation:
the PhysX TGS solver remains stable across the entire tested range
($E \in [12.6, 1002.6]$~MPa) at a practical timestep.  This means
domain randomisation over $E$ can be performed safely in Isaac Sim
without the need for per-sample timestep adaptation.

\subsection{Limitations}

\begin{itemize}[nosep]
    \item \textbf{No axial elasticity}: the rigid capsule links cannot stretch.
          For most manipulation tasks this is acceptable, but applications
          requiring accurate elongation under tension would need soft
          translational springs.
    \item \textbf{Sim-to-sim validation only}: the stability sweep compares
          against the Govoni paper's simulated results.  Physical validation
          against a real cable with known $E$ remains future work.
    \item \textbf{Computational cost}: a 200-link chain at 64 solver iterations
          is more expensive than lower-fidelity alternatives.
          Discretisation sweeps (varying $N$) would help identify the minimum
          link count for a given task.
    \item \textbf{Bending damping does not damp bulk modes}: the cable's
          pendulum-mode oscillations are dissipated only by translational drag,
          not by the material-derived joint damping.  For applications requiring
          accurate full-cable energy decay, an additional linear damping term
          tuned to match real-world swing decay would be needed.
\end{itemize}


% ------------------------------------------------------------------
\section{Proposed Next Steps}
% ------------------------------------------------------------------

\begin{enumerate}[nosep]
    \item \textbf{Robot integration}: attach the cable to a manipulator
          end-effector in Isaac Sim and define a cable-manipulation task
          (e.g., routing a cable through a set of clips).
    \item \textbf{Reinforcement learning}: leverage domain randomisation
          over material parameters ($E$, $\zeta$, cable length) to train
          robust manipulation policies.  My stability results confirm that
          randomising $E$ up to \SI{1}{GPa} will not produce divergent
          training samples at the standard Isaac Sim timestep --- a concern
          raised explicitly by Govoni et al.
    \item \textbf{Physical validation}: measure $E$ for a real cable specimen
          and compare simulated bending profiles against physical measurements.
    \item \textbf{Discretisation study}: sweep $N \in \{10, 50, 100, 200\}$
          to quantify the trade-off between visual fidelity, physical accuracy,
          and computational cost.
\end{enumerate}


% ------------------------------------------------------------------
% References
% ------------------------------------------------------------------
\begin{thebibliography}{9}
\bibitem{govoni2025}
    A.~Govoni, N.~Zubair, S.~Soprani, and G.~Palli,
    ``Performance Analysis of a Mass-Spring-Damper Deformable Linear Object
    Model in Robotic Simulation Frameworks,''
    \textit{arXiv preprint arXiv:2504.13659}, 2025.

\bibitem{lv2017}
    N.~Lv et al.,
    ``Physically based real-time interactive assembly simulation of cable
    harness,''
    \textit{Journal of Manufacturing Systems}, vol.~43, pp.~385--399, 2017.

\bibitem{bergou2008}
    M.~Bergou, M.~Wardetzky, S.~Robinson, B.~Audoly, and E.~Grinspun,
    ``Discrete Elastic Rods,''
    \textit{ACM Trans.\ on Graphics (SIGGRAPH)}, vol.~27, no.~3, 2008.

\bibitem{bhasin2006}
    Y.~Bhasin and A.~Liu,
    ``Bounds for damping that guarantee stability in mass-spring systems,''
    \textit{Studies in Health Technology and Informatics}, vol.~119,
    pp.~55--60, 2006.
\end{thebibliography}

\end{document}
